%%% MA_04_Economic_Policy.tex -- reorganised; frames moved from the old MA_xx files (see comments)

\section{Economic Policy
% : Rationale, Institutions, Objectives, Instruments
}


\begin{frame}{Micro- vs. Macroeconomics}
    \mycomment{
     Macroeconomics: where Microeconomics looks at the outcomes of individual markets, macroeconomics looks at the outcomes of all markets simultaneously.

     Thus, it looks at larger aggregates. For example: rather than the market of one particular good, the market for a class of goods, or for goods in general.
    
     Macroeconomics is thus concerned with overall equilibrium, growth, and the state of the economy as a whole, including distribution of income and wealth.

     Micro- and Macroeconomics are not truly \enquote{branches} of Economics, but rather different approaches. Many topics, such as taxes, can be analysed using microeconomic methodology. 
    }
    \showmytree[]{../shared/tikz/forest-micromacro}
\end{frame}

%%% [moved from MA_0_Organisation_and_Module_Walkthrough.tex, frame 2]
\begin{frame}{Micro- vs. Macroeconomics}
    \mycomment{
    Microeconomics isolates individual markets from the rest of the economy and looks at effects that changes to these markets have on the equilibrium of that market or those selected markets. Thus, it deals with partial equilibria.    

    For example consider savings: if you cut your spending, you can improve your savings. 
    }
    \showmytree[1]{../shared/tikz/forest-micromacro}
\end{frame}

%%% [moved from MA_0_Organisation_and_Module_Walkthrough.tex, frame 3]
\begin{frame}{Micro- vs. Macroeconomics}
    \mycomment{
    Macroeconomics looks at the entire economy, so at the simultaneous equilibrium on \enquote{all} markets. Thus, it deals with the general equilibrium.

    Now think again: if everyone cuts their spending, the result could be an economic downturn; nobody will be able to save more than before (paradox of thrift).

    Many things that are true for individual markets, and strategies that may work on individual markets, may not work for the economy as a whole.    

    Some say Macroeconomics should be firmly founded on microeconomic principles. Some hold that the whole is more than its parts and a perfect micro-foundation is not possible, or not sufficient.
    }
    \showmytree[2]{../shared/tikz/forest-micromacro}
\end{frame}

%%% [moved from MA_0_Organisation_and_Module_Walkthrough.tex, frame 4]
\begin{frame}{Economic Policy}
    \mytext{
    Economic policy asks how the state can influence the economy. 
    
    The influence on individual markets is a part of Microeconomics, or microeconomic policy. 
    
    Here, we are primarily interested in macroeconomic policy, or macroeconomic effects of microeconomic policy.

    \begin{itemize}
        \item What means are available to the state? What are the constraints?
        \item Apart from the influence on individual markets, what are the impacts (or side effects) of a single policy on the economy as a whole?
    \end{itemize}
    }
    \mycomment{
        Much can be said about the exact delineation between micro- and macroeconomics, the micro-foundation of macroeconomics, or if and how both are to some degree independent. Here, we have no room to dive into this academic debate. 

        Large non-governmental actors can of course also have an economic policy, or advise states to pursue such a policy (lobbying).
    }
\end{frame}



\begin{frame}{Shell Model}
    \mycomment{
    There are many levels of government, from the central government over federal states down to the communal level; and on each level, there are multiple ministries, government owned companies, and institutions.

    These sometimes complement, and sometimes compete. It is worthwhile and not an easy task to take a more differentiated view. In particular, government statistics need to be carefully read, for example on budget and debt level of government.
    }
    \showmytikz{../shared/tikz/tikz-shellmodel}
\end{frame}



\begin{frame}{Levels of Government}
    \showmytikz[1,2,3,4,5]{../shared/tikz/tikz-subsidiarity}
    \mycomment{
    Depending on the size of a state, it comprises of different sub-divisions. 
    
    Names, level of autonomy, and the responsibilities, of these units vary considerably between states.
    
    The level of autonomy, and the responsibilities, of these levels vary considerably between states. There are also sometimes supra-national levels, as the European Union.
    }
\end{frame}

\begin{frame}{Principle of Subsidiarity}
    \mytext{
    The smallest (lowest) government level that can tackle a problem should be in charge of this problem.
    }
    \mycomment{
    This (German) principle can also be found in European law (Art. 5(3) TEU).
    }
\end{frame}

\begin{frame}{Principle of Connexity}
    \mycomment{
    It happens that central government makes decisions that binds lower levels of government to provide, which incurs costs for the lover government level without funding provided from the higher level.

    Often much of the expenses on the lover level are obligatory, so that the lower levels have less discretion to decide what to do with their funds. 

    Thus, each (federal) law should be checked upon its budgetary consequences for all government levels. This is the connexity principle (German: \enquote{Konnexitätsprinzip}).
    }
    \mytext{
    The level of government whose decision incurs costs should also cover those costs. \enquote{Whoever orders, pays}.
    }
\end{frame}


\begin{frame}{Principle of Proportionality}
    \mycomment{
    The principle of proportionality is a concept from German (and Catholic Church) law, but generally applicable.

    For example: expropriation may solve a problem; but order policy (regulation) might solve the problem with less severe consequences.
    }
    \mytext{
    Every action of the state should be---
    \begin{enumerate}
        \item legitimate $\leadsto$ (the state has the right to do it),
        \item necessary (the state has reason to do it),
        \item suitable (the action has the desired effect),
        \item appropriate (there is no milder way to achieve the goal).
    \end{enumerate}
    }
\end{frame}

%%% [moved from MA_8_Economic_Policy.tex, frame 11]
\begin{frame}{Order vs. Process Policy}
    \mycomment{
    Economic policy can take a multitude of forms, which can be classified in a number of different ways. No one scheme is exhaustive.

    One major distinction is that between order policy, that is the government setting the framework for liberal economic activity, and process policy, which influences what is happening on markets directly. This is a specific German distinction (\noun{Eucken} \cite{Eucken1952}).
    }
    \showmytree[1]{../shared/tikz/forest-economic-policy}
\end{frame}

%%% [moved from MA_8_Economic_Policy.tex, frame 12]
\begin{frame}{Fiscal Policy}
    \mycomment{
    Within process policy, any policy that touches the government budget is fiscal policy. This includes the income side of the government: tax policy; and the expenditure side. 
    
    The government has a plethora of tools at its disposal. These can be more broadly applied, or more specific and discretionary. The government policy itself can be more or less rule based, depending on the constitutional and legal situation.
    }
    \showmytree[only 2]{../shared/tikz/forest-economic-policy}
\end{frame}

%%% [moved from MA_8_Economic_Policy.tex, frame 13]
\begin{frame}{Further Means of Economic Policy}
    \mycomment{
    Apart from fiscal policy, the government has other means to influence economic activities. One important branch is monetary policy, which resides with the Central Bank. This we will cover extensively later.

    Structural policy overlaps with fiscal policy; apart from the immediate fiscal effects, policies have spatial and structural effects. One important tool is the establishment of special economic areas, where incentives to invest are set by tax rebates, lower standards, or subsidies.

    The government can also promote industries or investment in certain areas or sectors by providing information services. Trade agreements and foreign policy also impact economic activities.
    }
    \showmytree[3]{../shared/tikz/forest-economic-policy}
\end{frame}

%%% [moved from MA_8_Economic_Policy.tex, frame 14]
\begin{frame}{Further Means of Economic Policy}
    \mycomment{


    Structural policy overlaps with fiscal policy; apart from the immediate fiscal effects, policies have spacial and structural effects. One important tool is the establishment of special economic areas, where incentives to invest are set by tax rebates, lower standards, or subsidies.

    The government can also promote industries or investment in certain areas or sectors by providing information services. And trade agreements and foreign policy also impacts economic activities.
    }
    \showmytree[only 4]{../shared/tikz/forest-economic-policy}
\end{frame}


%%% [moved from MA_8_Economic_Policy.tex, frame 15]
\begin{frame}{Why Economic Policy?}
    \mytext{
        We know from Microeconomics that competitive equilibria are Pareto-efficient.
        \begin{itemize}
            \item Why \textit{should} a government pursue economic policy?
            \item[$\rightarrow$] normative
            \item Why \textit{would} a government pursue economic policy?
            \item[$\rightarrow$] descriptive
        \end{itemize}        
        Economic Policy is intervention in an otherwise unruly world. It is necessary to establish markets; but it can do much more.
    }
\end{frame}

\begin{frame}{Arguments for Intervention}
    \mycomment{
    The discussion for and against having a state, what the state should do and what it should not do, is old and large. Microeconomic theory proofs that competitive markets are Pareto-efficient. From this follows that the state should primarily work to support such competitive markets. 
    Would a situation without state really be a war of all against everyone (\enquote{bellum omnium contra omnes}, as Thomas \noun{Hobbes} famously claimed?
    
    But the state can do much more. Firstly, it can correct individual markets where they fail. But then the sum of all state activities has an influence on the economy \textit{as a whole}. There are thus microeconomic policies---policies addressing individual markets---and macroeconomic policies, that steer the economy \textit{as such.}
    }
    \mytext{
        \begin{itemize}
            \item Correct market failure
            \begin{itemize}
                \item under-use of resources
                \item inequality (in chances, endowment)
                \item over-use of resources
                \item externalities, tragedy of the commons
            \end{itemize}
            \item The government may be wiser 
            \begin{itemize}
                \item than individuals (paternalism);
                \item than the markets
            \end{itemize}
            \begin{itemize}
                \item long-term strategy
                \item more information than subjects
                \item morally superior 
            \end{itemize}
            \item government can raise revenue
        \end{itemize}
    }
\end{frame}

\begin{frame}{Arguments for Intervention}
    \mytext{
    Correction of market failures
    \begin{itemize}
        \item adverse selection leading to lack of  insurance
        \begin{itemize}
            \item health insurance
            \item unemployment insurance
            \item pensions (old age insurance)
            \item ...
        \end{itemize}
        \item Monopolies
        \item Poverty traps
        \item Public goods
        \item Infrastructure
        \item Education
        \item ...
    \end{itemize}
    }
    \mycomment{
    The debate about scope and reason for government's intervention into markets, or economic policy, is vibrant. Maria \noun{Mazzucato} \cite{mazzucato2011entrepreneurial} 
    stresses the importance of government's entrepreneurial funding of education and research in particular; Elionor \noun{Ostrom} \cite{Ostrom1990} the importance of institutions for public goods and the commons. 
    }    
\end{frame}


%%% [moved from MA_8_Economic_Policy.tex, frame 17]
\begin{frame}{Possible Government Failures}
    \mylonglist{
    \begin{itemize}
        \item government not wiser
        \item populism leads to irrational policy
        \item chauvinism, favouritism, nepotism
        \item tends to be conservative rather than innovative (contested)
        \item self-interest of government
        \item interest of voters: median voter theorem
        \item excess burden, crowding out, deadweight loss
        \item price controls and nationalisation lead to over-use, low prices, little investment
        \item rent-extraction, corruption
        \item debt financing burdens future generations
        \item unsustainable debt: frequent crises
        \item ...
    \end{itemize}
    }
\end{frame}

 
\begin{frame}{Macroeconomic Policy}
    \mycomment{
    While microeconomics looks at individual, isolated markets, or sometimes at two or three related markets, macroeconomics looks at the sum (or better: product, or even: factorial; the interplay) of all markets. 

    So if an economic policy is large, it has macroeconomic effects. Since the state is large (in modern, developed, industrialised states), its policy has macro effects.
    }
    \mytext{
    The macroeconomic perspective adds to that:
    \begin{itemize}
        \item an invention in one market has an influence on (all) other markets
        \item again, there are income and substitution effects
        \item thus, there are macroeconomic effects: every policy (of a certain size) affects interest rates, growth, employment, etc.
        \item the effects of a policy on the system are complex; side effects, feedback loops, and budget constraints can have \enquote{surprising} effects.
        \item there can be policy measures targeting the macro level specifically. For instance, monetary policy affects all markets.
    \end{itemize}
    }   
\end{frame}

\begin{frame}{Macroeconomic Policy}
    \mycomment{
    Individual markets may be incomplete or fail, so that resources are allocated to activities and distributed in shares that are not in accordance to what society wants (\enquote{volonté generale}, Jean-Jaques \noun{Rousseau}). Thus just like macroeconomic policy corrects shortcomings in single markets, macroeconomic policy corrects the shortcomings of the interaction of all markets in an economy.
    }
    \mytext{
    If an economy where the state provides rules, guarantees property rights and enforces claims is not perfect, then we need     
    macroeconomic policy
    \begin{itemize}
        \item to stabilise the economy
        \item to create fairness and equity in distribution
        \item to allocate to national public goods
        \item to facilitate income / income growth
    \end{itemize}
    }
\end{frame}
 
\begin{frame}{Macroeconomic Policy}
    \mytext{

    \begin{itemize}
        \item micro-policy (into individual markets) is different from macro-policy 
        \begin{itemize}
            \item spending on one market: micro-policy
            \item government spending as such: macro-policy
        \end{itemize}
        \item micro- and macro-policy may conflict
        \item macro-policy targets chief macroeconomic targets, at the core: GDP and growth
        \item GDP and growth relate to employment, investment, distribution etc.
        \item these again depend on price level (inflation), that depends on money supply, exchange rate, etc.
    \end{itemize}
    }

\end{frame}


%%% [moved from MA_8_Economic_Policy.tex, frame 1]
\begin{frame}{Objectives of Macroeconomic Policy}
    \showmytikz[label=fig:magicsquare,step={1}]{../shared/tikz/tikz-economicsmagicsquare}
    \mycomment{
    The so-called \enquote{magic square} (\textit{Magisches Viereck}) is a popular design in the German discourse. This is commonly understood to be \enquote{gesamtwirtschaftliches Gleichgewicht} (macroeconomic balance) as required by law in  GG~Art.~109 and Par.~1~StabG.

    \foreignblockcquote{ngerman}{stabg1967}{
        Bund und Länder haben bei ihren wirtschafts- und finanzpolitischen Maßnahmen die Erfordernisse des gesamtwirtschaftlichen Gleichgewichts zu beachten. Die Maßnahmen sind so zu treffen, daß sie im Rahmen der marktwirtschaftlichen Ordnung gleichzeitig zur Stabilität des Preisniveaus, zu einem hohen Beschäftigungsstand und außenwirtschaftlichem Gleichgewicht bei stetigem und angemessenem Wirtschaftswachstum beitragen.
    }

    In the English literature, this is sometimes called \enquote{Kaldor's Diamond}, though \noun{Kaldor}'s publication \cite{Kaldor1971} came after the German legislation.

    In addition, the \enquote{Schuldenbremse} (German balanced budget amendment) limits the yearly budget deficit to 0,35\,\% of GDP.

      }
\end{frame}


%%% have something about the budget constraint


\begin{frame}{Objectives of Macroeconomic Policy}
    \showmytikz[2]{../shared/tikz/tikz-economicsmagicsquare}
    \mycomment{
    There have been proposals to amend the \enquote{magic square} / \enquote{Kaldor's diamond} \myfigref{fig:magicsquare}. Sustainability should be taken into account \cite{Dullien2015Realitaetscheck}. However, sustainability is difficult to measure, and a long-term constraint rather than a short-term target. GDP, inflation etc.\ are indicators we have, and which can be influenced directly and quickly.

    In any case, it remains questionable if, or to which extent, these goals compete with one another, if they can indeed only be achieved in unison, or if one or the other should have priority.
    }
\end{frame}

%%% [moved from MA_8_Economic_Policy.tex, frame 3]
\begin{frame}{Objectives of Macroeconomic Policy - EU}
    \showmytikz{../shared/tikz/tikz-119TFEU}
    \mycomment{There is a similar concept in Art. 119 TFEU on the EU level. 

        \blockcquote[TFEU Art. 119 III]{tfeu_2012}{
            These activities of the Member States and the Union shall entail compliance with the following guiding principles: stable prices, sound public finances and monetary conditions and a sustainable balance of payments.
        }
    
    }
\end{frame}

%%% [moved from MA_8_Economic_Policy.tex, frame 4]
\begin{frame}{Objectives of Macroeconomic Policy - EU}
    \showmytikz{../shared/tikz/tikz-3TEU}
    \mycomment{
    Art. 3 TEU defines the basis and the aims of the internal market of the Union.

    \blockcquote[TEU Art. 3,3]{TEU2012}{
        The Union shall establish an internal market. It shall work for the sustainable development of Europe based on balanced economic growth and price stability, a highly competitive social market economy, aiming at full employment and social progress, and a high level of protection and improvement of the quality of the environment. It shall promote scientific and technological advance.    
    }
    
    }
\end{frame}

\begin{frame}{Euro Convergence Criteria}
    \mytext{
    To enter the third stage of the Economic and Monetary Union and adopt the Euro, states are obliged to fulfil five criteria. The basis is Art. 140 and 162 TFEU, but there are additional laws.

    \begin{itemize}
        \item HICP inflation below an average (actual formula is complex)
        \item yearly government budget deficit per GDP $<3\,\%$ 
        \item government debt per \mbox{GDP $<60\,\%$}
        \item exchange rate stability: participation in the European Monetary System (EMS)
        \item long-term interest rates (average yields for 10-year government bonds in the last year) not higher than 2 percentage point than the 3 EU member states with lowest HICP inflation
    \end{itemize}
    }
    \mycomment{
    HICP is the Harmonised Index of Consumer Prices, an indicator published by the ECB.
    }
\end{frame}

\begin{frame}{Euro Convergence Criteria}
    \mycomment{
    These so-called \enquote{Maastricht Criteria} resemble the \enquote{magic square} of German policy (\enquote{Kaldor's Diamond}---but instead of full employment, there is the \enquote{fiscal criterion}. This comes from the fact that a monetary union can only function if no state is \enquote{free riding}. There are no \enquote{EU bonds}; states issue their own bonds and manage their own debt. If this is done excessively, that  could endanger the monetary union or force some states to bail out others.
    }
    \showmytikz[1]{../shared/tikz/tikz-euroconvergencecriteria}
\end{frame}



\begin{frame}{Musgrave's Three Functions}
    \showmytikz{../shared/tikz/tikz-musgrave}
    \mycomment{
    Richard \noun{Musgrave} established the distinction between three  policy objectives, and a way to classify economic policy \cite{Musgrave1959} .
    \begin{itemize}
        \item allocation (vs. market failures): to what should resources be used? The outcome of free markets may be inefficient, as they may not provide or underprovde some goods (public goods); they may also over-use resources (commons). 
        \item distribution (equity): some areas may be underdeveloped; people may need help (the young, the elderly, the sick); distribution of income, wealth, or chances may be considered unfair and need correction only the state can provide. 
        \item stabilisation (business cycles): the macroeconomic development of an economy may be prone to (wild) swings, or even crisis, and require steering. This is macroeconic policy.
    \end{itemize}
    }
\end{frame}

\begin{frame}{Musgrave's Three Functions}
    \showmytikz[1]{../shared/tikz/tikz-musgrave}
    \mycomment{
        The three functions (and corresponding policies) may not be independent. More efficient markets may influence equity, etc.

        Macroeconomic policy is dealing with stabilisation policy. Nevertheless, allocation and distribution policy may have macroeconomic consequences.

        A policy to change distribution may or may not be efficient; and an allocation policy may have influences on distribution. So generally, every policy should be judged by looking at all three aspects.
    }
\end{frame}


%%% [moved from MA_8_Economic_Policy.tex, frame 9]
\begin{frame}{Market Conformity of Economic Policy}
    \mytext{
    Economic Policy should take the insights from Microeconomics into account:
    \begin{itemize}
        \item prices depend on volumes, and volumes on prices
        \item dead weight loss / excess burden
        \item crowding out
        \item income and substitution effects
        \item ...
    \end{itemize}
    Market conformity: rules that work \textit{through} the price mechanism (such as taxes and subsidies) rather than suspending it (such as price controls) \cite{Eucken1952}.
    }
    \mycomment{
    The concept of \enquote{market conformity} is a core concept of German \enquote{Soziale Marktwirtschaft} and associated with \noun{Röpcke} \cite{roepke1942gesellschaftskrisis}, \noun{Müller-Armack} \cite{muellerarmack1947wirtschaftslenkung} and \noun{Eucken} \cite{Eucken1952}. 
    
    There is critique of even this type of intervention. On is the \enquote{vicious cycle  of intervention}: every intervention has side effects that lead to more intervention \cite{vonmises1929kritik}. 
    }
\end{frame}


\begin{frame}{Tinbergen  Rule}
    \mytext{
    To achieve a policy goal, you need a lever.

    \fbox{
    \colorbox{yellow!40}{
    \parbox{0.8\textwidth}{\centering
        To hit $n$ independent targets, you need\\ at least $n$ independent instruments.
        }
    }
    }
    
    This is known as \noun{Tinbergen} Rule \cite{tinbergen1952theory}. It is a fundamental principle of economic policy design.

    Questions to ask--:

    \begin{itemize}
        \item How many tools does the state actually have?
        \item What are the restrictions the state is facing?
        \item What are the unwanted side effects of a policy?
        \item What are the unwanted side effects of a set of policies?
        \item Are these levers truly independent?
    \end{itemize}
    }

\end{frame}

\begin{frame}{Macroeconomic Targets and Policy Levers}
    \mytext{
    Often, the government is clearer about objectives than about tools (\enquote{diminish poverty} - but how?). 
    
    Which tools should be used for which target? The one that fits best, obviously.

    \begin{itemize}
        \item assign each instrument to one target
        \item assign to objective that has the comparative advantage (\enquote{\noun{Mundell}'s Assignment Principle})
    \end{itemize}
    Which tool is suited for which target is one of the chief questions macroeconomic theory is set to answer. 
    }
    
    \mycomment{
    It would be too simple to think that one measure could heal all problems. For instance, nationalisation of monopolies may allow the new owner (the state) to lower the price. But another effect will not have been solved: the new monopoly (the state monopoly) may offer the goods for lower prices, but still underprovide or underinvest. 

    Similarly, lower interest rates can improve employment, but will also increase the price level, which might even offset that effect. It might turn out not to be the right tool for this target altogether.
    
    If a policy has just one target, the success could be measured. If a policy has more than one target, measuring the success is difficult. And if multiple policies have multiple targets and contradicting side effects, the policy mix might create more chaos than heal existing problems.
    }
\end{frame}


\begin{frame}{Macroeconomic Targets and Policy Levers}
    \mycomment{
    \noun{Mundell} came to a different conclusion; today, the assignment is typically:

    \begin{itemize}
        \item monetary policy --  stabilisation
        \item fiscal policy -- redistribution and allocation
        \item[] (and stabilisation)
        \item order policy  -- allocation 
        \item[] (and also distribution)
    \end{itemize}

    That is, of course, not the end of the story. In particular, fiscal policy has a major influence on macroeeconomic stability, as the government can extend its expenditure to stimulate the economy, and incur government debt, which in the long run might destabilise the economy. 
    
    The interplay of fiscal and monetary policy is an important field of study in macroeconomics.
    }
    \showmytikz[2]{../shared/tikz/tikz-musgrave}
\end{frame}


\begin{frame}{Macroeconomic Targets and Policy Levers}
    \mytable{
      \documentclass{standalone}
\usepackage{hsbcommon}
\usepackage{array}

%% Show/hide steps for standalone work — edit this line:
\makeatletter\AtBeginDocument{\@buildstepset{1,2,3,4,5,6,7}}\makeatother

\begin{document}
\newcolumntype{L}[1]{>{\raggedright\arraybackslash}p{#1}|}%
\newcolumntype{T}{L{2.6cm}L{2.6cm}L{2.6cm}L{2.6cm}}%
    \begin{tabular}{T}
        \hline
         & Allocation & \colorbox{blue!20}{Distribution} & \colorbox{red!20}{Stabilisation} \\
        \hline
        Order policy
         & \cellcolor{green!10}\textbf{primary}\par property rights,\par competition law
         & \textbf{secondary}\par labour law, inheritance law
         & framework:\par monetary and fiscal rules \\
        \hline
        \colorbox{blue!40}{Fiscal policy}
         & \textbf{secondary}\par\par public goods, Pigouvian taxes
         &\cellcolor{blue!10}\textbf{primary}\par taxes, transfers
         & \textbf{secondary;}\par primary in the aggregate \\
        \hline
        \colorbox{red!40}{Monetary policy}
         & ---
         & side effects (asset prices)
         & \cellcolor{red!10}\textbf{primary} \\
        \hline
        \end{tabular} 
\end{document}
    }
\end{frame}



\begin{frame}{Government Budget: Primary and Total}
    \mytext{
    The government faces a budget constraint.
    \begin{equation}
        \begin{split}            
          & \quad \text{expenditure } G_t\\
        - &\quad  \text{revenue }T_t \\
        = &\quad  \textbf{primary deficit }  
        \end{split}
    \end{equation}

    If you add the spending of interest on the outstanding debt, you get to the \textbf{total deficit}:
    \begin{equation}
        \begin{split}            
        & \quad \text{expenditure } G_t\\
        + &\quad  \text{interest } iB_{t-1}\\
        - &\quad  \text{revenue }T_t \\
         = &\quad  \textbf{total deficit } D_t\\
         = &\quad  B_t - B_{t-1}
        \end{split}
    \end{equation}
    }
    \mycomment{
        The Maastricht criterion refers to the total deficit.
    }

\end{frame}

\begin{frame}{Government Budget: Example}
    \showmytikz[example]{../shared/tikz/tikz-budgetbars}
\end{frame}

\begin{frame}{Government Bonds}
    \mytext{
    Government issue debt in a variety of forms, most commonly by issuing bonds: 
    
    \begin{itemize}
        \item Treasuries~(US)
        \item gilts~(UK)
        \item Bundesanleihen (\enquote{Bunds} in English) (Germany)
        \item Landesschatzanweisungen (German Länder)
        \item OATs~(France)
        \item JGBs~(Japan)
    \end{itemize}
    }    
    \mycomment{
        Investors are closely watching the \textbf{spread} between those treasury bonds, for example the spread between German and French bonds, both in Euro.    
    }
\end{frame}

\begin{frame}{Government Bonds}
    \showmytikz{../shared/tikz/tikz-bonds-yields}
    \mycomment{
        Depending on the actual bill, there may be an annual (or twice annual) payouts (per \textbf{coupon}). Then there is a difference between \textbf{Yield to Maturity YTM}  and \textbf{Current Yield} (simple annual returns), not shown in this illustration.

        The debt issued can have long maturity (10 years or more) or short maturity. The spread between the yield of the long and short term loans is telling (\textbf{yield curve}). Normally, long term bonds \textit{normally }yield higher interest rates than short term bonds.     

        Government bonds are highly standardised bonds, and easy to sell and re-sell. Indeed, instead of holding one million US dollar in cash, you'd typically buy US government bonds with that cash as a \enquote{safe haven}.
   
    }
\end{frame}

\begin{frame}{Government Bonds}
    \showmytikz[1,2]{../shared/tikz/tikz-bonds-yields}
    \mycomment{
        Bonds are sold at auction or on exchanges at a price below their face value; the result is a \textbf{yield} (interest).
        
        \textbf{Yields move in the opposite direction of prices.}

        Because bonds have different denominations and coupons, the one interesting figure to compare them all is the yield.

        Bonds compete with other assets over investor's money. So yields typically follow (expected, real) interest rates in the economy---plus a risk premium.
    }
\end{frame}

\begin{frame}{Government Bonds}
    \mytext{
    The interest rate is typically lower than private borrowing:

    \begin{itemize}
        \item states never die
        \item the entire economy (GPD) is collateral, as the state can raise taxes
        \item the markets for government bonds are huge, the bonds highly liquid assets
        \item USA: \enquote{exorbitant privilege}: reserve currency / currency in which prices are quoted/invoiced
    \end{itemize}
    }
    \mycomment{    
        States can borrow as long as investors trust them to repay, but states can and do default.

        Historical observation from the Euro crisis 2011-2012: interest rates on government bonds rise above 6\,\%, markets seem to mistrust the ability of the state to repay its debt.
    }    
\end{frame}

% \begin{frame}{Budget Data for Comparison}
    \mytext{
    The following are illustrations showing budgets of central government, social security systems, and the city state of Bremen. 

    France’s welfare state is large relative to its economy, not unusually generous per person. Germany spends about as much per head, but from a larger GDP per head.
    }
    \mycomment{
        If you were to compare these across each other, you must account for the following:
        \begin{itemize}
            \item differences in population: Germany has roundabout 25\,\% more population than France
            \item differences in GDP per head: a richer country will have more income and can afford more expenditure per head than a poorer one.
            \item purchasing power parity: the pure numbers all are denominated in EUR, but need to be adjusted to be meaningful.
        \end{itemize}
    }
\end{frame}

\begin{frame}{Germany's State Budget (2027 draft) - Federal}
    \mycomment{
    The by far largest position in the federal budget of Germany is \enquote{Arbeit und Soziales} (BMAS), followed by defence (BMVg). The third post is \enquote{Allgemeine Finanzverwaltung}; debt service comes fourth, followed by traffic (here BMV). All other positions are smaller; the ministry for economic affairs (BMW) only receives about 1\,\% of the federal budget.

    Social security, unemployment etc.\ is only partly included here: most of it runs via contributions and not through the State budget. \enquote{Social} includes that part that does (\enquote{Grundsicherung}; subsidies to pensions, etc.).

    On the income side, about three quarters are taxes, levies etc., with another fifth \enquote{Bundesschuld}, new debt. Both sides are equal, in 2027 roundabout 555 billion Euro.
    }
    \showmytikz[1]{../shared//tikz/tikz-budgetbars}
\end{frame}

\begin{frame}{Germany's State Budget (2026) - Bremen}
    \mycomment{
    Germany is a federal state where some areas are exclusive priviledge of the Länder (the states). Most notably education, culture, and policing. So these positions do not appear on the federal level, but on the Länder level.

    According to the German constitution, tax revenue is shared between Bund, Länder and Kommunen (apart from some taxes that go directly to either Bund or Länder). As a result, the sum of all Länder budgets is be roughly the same as the federal budget. The federal state, however, borrows more, mostly outside of core budget (\enquote{Sondervermögen}). Also, substantial sums are moved between the different levels, so the budgets cannot simply be summed up.

    Bremen has the highest debt level of all German Länder by a margin, at approximately 37.000 Euro per inhabitant. (Federal debt adds another 22.000 Euro per head.) Bremen's debt is high relative to GDP, but its interest burden is moderate for now, thanks to low rates locked in during the past decade and federal consolidation aid. As old loans are refinanced at higher rates, the debt service segment will grow.
    }
    \showmytikz[2]{../shared//tikz/tikz-budgetbars}
\end{frame}

\begin{frame}{France: State Budget}
    \mycomment{For comparison, we can look at France. Of course, the numbers are not exactly comparable, since statistic methods may vary. But the picture is clear. The central government pays the teachers. Again, social insurance (pensions, health)  is mainly outside (contributions to a system independent of the state budget). However, the size of debt service and the size of new debt per year is striking.
    }
    \showmytikz[3]{../shared//tikz/tikz-budgetbars}
\end{frame}

\begin{frame}{Germany Social Insurance Budget}
    \mycomment{
    Germany's social security system is running deficits which are covered by federal funds, so this appears in the federal budget's position \enquote{social}. Unemployment scheme is among \enquote{div.}.


    }
    \showmytikz[4]{../shared//tikz/tikz-budgetbars}
\end{frame}

\begin{frame}{France Social Insurance Budget}
    \mycomment{
    In France, the social insurance system issues their own debt, and runs its own deficit. Apart from contributions, there is a general tax \enquote{contribution sociale généralisée} (CSG) on almost all income to finance the system, which could indeed be compared to a state subsidy.

    Some items still sit outside, most notably France's mandatory supplementary pensions (Agirc-Arrco, roughly 100 bio EUR) and unemployment insurance (Unédic, roughly 40 bio EUR).

    Note the central figures: the social security budget of France is larger than the state's budget. 
    }
    \showmytikz[5]{../shared//tikz/tikz-budgetbars}
    
\end{frame}

\begin{frame}{Government Spending to GDP Ratio}
    \mytext{
    Just how big is the state? To answer this, we measure government expenditure, including interest payments, relative to GDP.
 
    In 2025, the figures were 
    \begin{itemize}
        \item 50.3\,\% for Germany and
        \item 57.3\,\% for France.
    \end{itemize}
    The size of the French one comes chiefly from the larger social security budget.

    In contrast, Nigeria's ratio stands at about 12.5\,\% only.
    }
    \mycomment{
    \begin{equation}
        g_t=\frac{
          G_t  +   iB_{t-1}
        }{
        Y_t
        }
    \end{equation}    
    
    }
\end{frame}

\begin{frame}{Total Government Deficit to GDP Ratio}
    \mycomment{
    The ratio of the total deficit to GDP is:
    \begin{equation}
        d_t=\frac{G_t-T_t+iB_{t-1}}{Y_t}% = g_t-\tau_t 
    \end{equation}
    % with the revenue ratio $\tau_t=T_t/Y_t$.

    The Euro convergence criterium (Maastricht) is:
    \begin{equation}
        d_t\leq 3\,\%
    \end{equation}
    }
    \mytext{
    In 2025, this was 
    \begin{itemize}
        \item 2.7\,\% for Germany
        \item 5.1\,\% for France.
    \end{itemize}
    
    For 2026, Germany is expected at about 4\,\% and France at 5.4\,\%.
    }
\end{frame}

\begin{frame}{Government Debt Level to GDP Ratio}
    \mytext{
    At the end of 2025, it stood at 
    \begin{itemize}
        \item 63.5\,\% in Germany
        \item 115.7\,\% in France.
    \end{itemize}
    Both are rising.
    }
    \mycomment{
    
    Debt level to GDP:
    \begin{equation}
       b_t=\frac{B_t}{Y_t} 
    \end{equation}
    
    The Euro convergence criterion (Maastricht) is:
    \begin{equation}
        b_t\leq 60\,\%
    \end{equation}
    }
\end{frame}

\begin{frame}{Government Debt Sustainability}
    \mytext{
    Debt accumulates from deficits: $B_t=B_{t-1}+D_t$.
    
    Debt ratio is stable when debt grows as fast as GDP.
    
    Example: 
    \begin{itemize}
        \item GDP $100 \to 105$ ($\gamma=5\,\%$)
        \item debt $60 \to 63$ ($d=3\,\%$):
        \item the ratio stays at $63/105=60\,\%$.
    \end{itemize}
    
    This is the logic linking the two Maastricht criteria. Those interest, growth rates and debt levels were realistic in 1992.
     
    }
    \mycomment{
    This formula shows the connection between (stable) deficit $b^*$, growth rate $\gamma$ and deficit ratio $d$:
     \begin{equation}
            b^*=\frac{1+\gamma}{\gamma}\,d \approx \frac{d}{\gamma}
    \end{equation}
    It can be shown that at nominal growth of around 3\,\%, a 3\,\% deficit stabilises debt at about 100\,\%, not 60\,\%.

    The interest rate does not appear in this formula, because the total deficit $d$ already contains the interest payments. But the interest rate decides what is left for everything else: if $i>\gamma$, interest alone exceeds the deficit that keeps debt stable, so the government needs a primary surplus.
    
    % All of this still assuming that $\gamma>i$; if not, a stable debt level with a primary deficit is impossible.
    
    %     Dividing by $Y_t$, with $\gamma_t$ as nominal GDP growth
    %     (real growth plus inflation), and using $Y_{t-1}/Y_t=1/(1+\gamma_t)$:
    %     \begin{equation}
    %         b_t = \frac{b_{t-1}}{1+\gamma_t} + d_t
    %     \end{equation}
            
    %     The debt ratio is stable ($b_t=b_{t-1}=b$) when debt grows as fast as GDP, which gives
    %     \begin{equation}
    %         b=\frac{1+\gamma}{\gamma}\,d \approx \frac{d}{\gamma}
    %     \end{equation}

    % The total deficit contains interest payments. Splitting it into the     primary deficit $p_t=(G_t-T_t)/Y_t$ and interest gives:
    % \begin{equation}
    %     b_t = \frac{1+i}{1+\gamma_t}\,b_{t-1} + p_t
    % \end{equation}
    % \begin{equation}
    %     \quad\Rightarrow\quad
    %     b_t-b_{t-1} \approx (i-\gamma_t)\,b_{t-1} + p_t
    % \end{equation}
    % If the interest rate exceeds growth ($i>\gamma$), debt grows by itself, and a primary surplus is needed to stabilise it.
    % If growth exceeds the interest rate ($\gamma>i$), even a small primary deficit is sustainable.
    }
\end{frame}

\begin{frame}{Government Debt Sustainability}
    \mytable{
    For GDP 100, debt 100, growth $\gamma=3\,\%$:
    
    \begin{tabular}{llll}
        \toprule
        interest & interest    & total deficit & primary balance \\
        rate    &  paid       & allowed       & must be \\
        \midrule
        $i=2\,\%$ ($i<\gamma$) & 2 & 3 & deficit of 1 \\
        $i=3\,\%$ ($i=\gamma$) & 3 & 3 & balanced (0) \\
        $i=5\,\%$ ($i>\gamma$) & 5 & 3 & surplus of 2 \\
        \bottomrule
        \end{tabular}
    }   
\end{frame}

\begin{frame}{Debt Sustainability}
    \mycomment{
    Debt is not a problem when it helps generate a growth of income that allows to cover the interest due. This is true for any institution, or household. For the state, the source of income is GDP (that is taxed).

    The sustainability of debt depends on the ratio of the growth of income, the growth of debt (the interest rate) and the debt level.

    For our graphic example, we show a debt level $b_1$ of $50\,\%$ of income. To make the argument more visible, we work with very large rates for growth and interest (50\,\% and 25\,\%, respectively).

    If income grows twice as fast as debt, it is possible to take on more debt and retaining the same ratio. This additional debt equals the deficit, as it covers expenses above income.

    If this remains like this \textit{forever,} then it is \textit{forever }possible to take on more debt. You \enquote{grow out of debt}.
    }
    \showmytikz[1]{../shared/tikz/tikz-sustainabledebt}
\end{frame}

\begin{frame}{Debt Sustainability}
    \mycomment{
        Remember the law of diminishing returns? Trees don't grow into the sky. At some point, growth may slump, or interest rates may rise---or both. 
        The \enquote{normal} case is that GDP growth is \textit{lower} than the interest rate.
    
        If the interest rate is higher than the rise in income, the debt level will rise, and continue to do so, as long as the relation stays this way.

        Then it is necessary to \enquote{save some money elsewhere}, that is to have a surplus; this surplus must be used to reduce the debt, so that debt level remains stable, or goes down.

        Needless to say: to get from one regime into the other is painful.
    }
    \showmytikz[4]{../shared/tikz/tikz-sustainabledebt}
\end{frame}

\begin{frame}{Government Debt Sustainability}
    \mycomment{
    Assume a balanced primary budget ($p_t=0$: taxes cover all spending except interest) and a debt ratio of 100\,\%. Interest is then paid with new borrowing, so only the gap between the interest rate and nominal growth (real growth plus inflation) drives the path.
    
    Then:
    \begin{itemize}
        \item if $i>\gamma$, debt grows through interest faster than GDP grows, so the debt ratio explodes (unsustainable);
        \item if $i=\gamma$, debt and GDP grow at the same rate, so the debt ratio stays constant;
        \item if $i<\gamma$, GDP grows faster than debt, so the debt ratio shrinks, without any repayment.
    \end{itemize}
    Note that inflation lowers the debt ratio only as long as interest rates do not follow. In the long run, investors demand higher interest rates when inflation is higher, so $i$ rises along with $\gamma$ and the effect largely disappears. Only existing debt at fixed rates is permanently devalued.
    }
    \showmytikz{../shared/tikz/tikz-debt-interest-growth}

\end{frame} 

\begin{frame}{Government Debt Sustainability}
    \mytext{
    Government debt is sustainable only if---
   \begin{itemize}
       \item interest rate below GDP growth ($i<\gamma$): a moderate primary deficit is possible, or
       \item interest rate above GDP growth ($i>\gamma$): a sufficient primary surplus is needed,
       \item and investors keep trusting the state to repay.
   \end{itemize}
    }
\end{frame}

